\subsection{最初に必要な用語}
\paragraph{表1：この章で最初に押さえる用語}
\par\smallskip\noindent\refstepcounter{table}\label{tab:ch10-01}表 \thetable: 表1：この章で最初に押さえる用語における用語・初学者向けの説明\par\smallskip
表 \thetable は、表1：この章で最初に押さえる用語における用語・初学者向けの説明を比較・確認しやすい形で整理したものである。\par\smallskip
\begin{longtable}{>{\raggedright\arraybackslash}p{0.28\linewidth}>{\raggedright\arraybackslash}p{0.64\linewidth}}
\toprule
用語 & 初学者向けの説明\\
\midrule
\endfirsthead
\toprule
用語 & 初学者向けの説明\\
\midrule
\endhead
プロセス & 入力を受け取り、作業活動によって出力へ変換するまとまりである。ここでは人間の作業活動を指す。\\
活動 & プロセスを構成するまとまった作業である。例として設計、実装、検証がある。\\
タスク & 活動をさらに細かくした、実行すべき作業単位である。\\
\term{ライフサイクル}{Life Cycle} & 構想から廃止まで、システムや製品がたどる全体の変化である。\\
ライフサイクルモデル & 作業をどの順序で、どの程度反復し、どの時点で成果物を出すかを示す型である。\\
\term{プロジェクト}{Project} & 明確な開始と終了を持ち、固有の製品・サービス・結果を作る一時的な取り組みである。\\
測定 & 状態や成果を数値や指標で把握する活動である。プロセス改善の土台になる。\\
経験的証拠 & 実際の測定や観察から得た根拠である。プロセス選択や改善判断を支える。\\
\bottomrule
\end{longtable}
